% Options for packages loaded elsewhere
\PassOptionsToPackage{unicode}{hyperref}
\PassOptionsToPackage{hyphens}{url}
%
\documentclass[
  ignorenonframetext,
]{beamer}
\usepackage{pgfpages}
\setbeamertemplate{caption}[numbered]
\setbeamertemplate{caption label separator}{: }
\setbeamercolor{caption name}{fg=normal text.fg}
\beamertemplatenavigationsymbolsempty
% Prevent slide breaks in the middle of a paragraph
\widowpenalties 1 10000
\raggedbottom
\setbeamertemplate{part page}{
  \centering
  \begin{beamercolorbox}[sep=16pt,center]{part title}
    \usebeamerfont{part title}\insertpart\par
  \end{beamercolorbox}
}
\setbeamertemplate{section page}{
  \centering
  \begin{beamercolorbox}[sep=12pt,center]{part title}
    \usebeamerfont{section title}\insertsection\par
  \end{beamercolorbox}
}
\setbeamertemplate{subsection page}{
  \centering
  \begin{beamercolorbox}[sep=8pt,center]{part title}
    \usebeamerfont{subsection title}\insertsubsection\par
  \end{beamercolorbox}
}
\AtBeginPart{
  \frame{\partpage}
}
\AtBeginSection{
  \ifbibliography
  \else
    \frame{\sectionpage}
  \fi
}
\AtBeginSubsection{
  \frame{\subsectionpage}
}
\usepackage{amsmath,amssymb}
\usepackage{lmodern}
\usepackage{iftex}
\ifPDFTeX
  \usepackage[T1]{fontenc}
  \usepackage[utf8]{inputenc}
  \usepackage{textcomp} % provide euro and other symbols
\else % if luatex or xetex
  \usepackage{unicode-math}
  \defaultfontfeatures{Scale=MatchLowercase}
  \defaultfontfeatures[\rmfamily]{Ligatures=TeX,Scale=1}
\fi
\usetheme[]{Montpellier}
% Use upquote if available, for straight quotes in verbatim environments
\IfFileExists{upquote.sty}{\usepackage{upquote}}{}
\IfFileExists{microtype.sty}{% use microtype if available
  \usepackage[]{microtype}
  \UseMicrotypeSet[protrusion]{basicmath} % disable protrusion for tt fonts
}{}
\makeatletter
\@ifundefined{KOMAClassName}{% if non-KOMA class
  \IfFileExists{parskip.sty}{%
    \usepackage{parskip}
  }{% else
    \setlength{\parindent}{0pt}
    \setlength{\parskip}{6pt plus 2pt minus 1pt}}
}{% if KOMA class
  \KOMAoptions{parskip=half}}
\makeatother
\usepackage{xcolor}
\IfFileExists{xurl.sty}{\usepackage{xurl}}{} % add URL line breaks if available
\IfFileExists{bookmark.sty}{\usepackage{bookmark}}{\usepackage{hyperref}}
\hypersetup{
  hidelinks,
  pdfcreator={LaTeX via pandoc}}
\urlstyle{same} % disable monospaced font for URLs
\newif\ifbibliography
\usepackage{color}
\usepackage{fancyvrb}
\newcommand{\VerbBar}{|}
\newcommand{\VERB}{\Verb[commandchars=\\\{\}]}
\DefineVerbatimEnvironment{Highlighting}{Verbatim}{commandchars=\\\{\}}
% Add ',fontsize=\small' for more characters per line
\usepackage{framed}
\definecolor{shadecolor}{RGB}{248,248,248}
\newenvironment{Shaded}{\begin{snugshade}}{\end{snugshade}}
\newcommand{\AlertTok}[1]{\textcolor[rgb]{0.94,0.16,0.16}{#1}}
\newcommand{\AnnotationTok}[1]{\textcolor[rgb]{0.56,0.35,0.01}{\textbf{\textit{#1}}}}
\newcommand{\AttributeTok}[1]{\textcolor[rgb]{0.77,0.63,0.00}{#1}}
\newcommand{\BaseNTok}[1]{\textcolor[rgb]{0.00,0.00,0.81}{#1}}
\newcommand{\BuiltInTok}[1]{#1}
\newcommand{\CharTok}[1]{\textcolor[rgb]{0.31,0.60,0.02}{#1}}
\newcommand{\CommentTok}[1]{\textcolor[rgb]{0.56,0.35,0.01}{\textit{#1}}}
\newcommand{\CommentVarTok}[1]{\textcolor[rgb]{0.56,0.35,0.01}{\textbf{\textit{#1}}}}
\newcommand{\ConstantTok}[1]{\textcolor[rgb]{0.00,0.00,0.00}{#1}}
\newcommand{\ControlFlowTok}[1]{\textcolor[rgb]{0.13,0.29,0.53}{\textbf{#1}}}
\newcommand{\DataTypeTok}[1]{\textcolor[rgb]{0.13,0.29,0.53}{#1}}
\newcommand{\DecValTok}[1]{\textcolor[rgb]{0.00,0.00,0.81}{#1}}
\newcommand{\DocumentationTok}[1]{\textcolor[rgb]{0.56,0.35,0.01}{\textbf{\textit{#1}}}}
\newcommand{\ErrorTok}[1]{\textcolor[rgb]{0.64,0.00,0.00}{\textbf{#1}}}
\newcommand{\ExtensionTok}[1]{#1}
\newcommand{\FloatTok}[1]{\textcolor[rgb]{0.00,0.00,0.81}{#1}}
\newcommand{\FunctionTok}[1]{\textcolor[rgb]{0.00,0.00,0.00}{#1}}
\newcommand{\ImportTok}[1]{#1}
\newcommand{\InformationTok}[1]{\textcolor[rgb]{0.56,0.35,0.01}{\textbf{\textit{#1}}}}
\newcommand{\KeywordTok}[1]{\textcolor[rgb]{0.13,0.29,0.53}{\textbf{#1}}}
\newcommand{\NormalTok}[1]{#1}
\newcommand{\OperatorTok}[1]{\textcolor[rgb]{0.81,0.36,0.00}{\textbf{#1}}}
\newcommand{\OtherTok}[1]{\textcolor[rgb]{0.56,0.35,0.01}{#1}}
\newcommand{\PreprocessorTok}[1]{\textcolor[rgb]{0.56,0.35,0.01}{\textit{#1}}}
\newcommand{\RegionMarkerTok}[1]{#1}
\newcommand{\SpecialCharTok}[1]{\textcolor[rgb]{0.00,0.00,0.00}{#1}}
\newcommand{\SpecialStringTok}[1]{\textcolor[rgb]{0.31,0.60,0.02}{#1}}
\newcommand{\StringTok}[1]{\textcolor[rgb]{0.31,0.60,0.02}{#1}}
\newcommand{\VariableTok}[1]{\textcolor[rgb]{0.00,0.00,0.00}{#1}}
\newcommand{\VerbatimStringTok}[1]{\textcolor[rgb]{0.31,0.60,0.02}{#1}}
\newcommand{\WarningTok}[1]{\textcolor[rgb]{0.56,0.35,0.01}{\textbf{\textit{#1}}}}
\setlength{\emergencystretch}{3em} % prevent overfull lines
\providecommand{\tightlist}{%
  \setlength{\itemsep}{0pt}\setlength{\parskip}{0pt}}
\setcounter{secnumdepth}{-\maxdimen} % remove section numbering
%\documentclass[unicode,10pt]{beamer}% 'unicode'が必要
\usepackage{luatexja}% 日本語を使う場合は必須
%\usepackage[japanese]{babel}	
%\usefonttheme[onlymath]{serif}
%\usepackage[T1]{fontenc} %これを使うと、ギリシャ文字が出ない
%\usepackage{textcomp}
%\usepackage[scale=1.0]{tgheros} % Sans serif font
%\usepackage[scaled]{beramono} % Monospace font
%\usepackage{luatexja-otf}
%\renewcommand{\kanjifamilydefault}{\gtdefault}
%\usepackage[haranoaji]{luatexja-preset}

% スライド番号の表示 %%%%%%%%%%%%%%%%%%%
\makeatletter
\setbeamertemplate{footline}{%
  \leavevmode%
  \hbox{%
  \begin{beamercolorbox}[wd=.5\paperwidth,ht=2.25ex,dp=1ex,right]{date in head/foot}%
    \insertframenumber{} \hspace*{2ex} % new version without total frames
  \end{beamercolorbox}}%
  \vskip0pt%
}
\makeatother
% スライド番号の表示 %%%%%%%%%%%%%%%%%%%
\ifLuaTeX
  \usepackage{selnolig}  % disable illegal ligatures
\fi

\title{第1章: イントロダクション}
\subtitle{Elena Llaudet \& Kosuke Imai.\\
Data Analysis for Social Science: A Friendly and Practical
Introduction.}
\author{}
\date{\vspace{-2.5em}2026-03-09}

\begin{document}
\frame{\titlepage}

\hypertarget{ux672cux66f8ux306eux6982ux89b3}{%
\section{1.1 本書の概観}\label{ux672cux66f8ux306eux6982ux89b3}}

\begin{frame}{本書の目的と構成}
\protect\hypertarget{ux672cux66f8ux306eux76eeux7684ux3068ux69cbux6210}{}
\begin{itemize}
\tightlist
\item
  社会科学におけるデータ分析の基本を、予備知識なしの状態から学ぶ。
\item
  統計ソフトRを用いて現実世界のデータを分析する。
\item
  計量社会科学研究の3つの基本的な目標を学ぶ：

  \begin{enumerate}
  \tightlist
  \item
    \textbf{測定} (Measurement)
  \item
    \textbf{予測} (Prediction)
  \item
    \textbf{説明}（因果推論、Causality）
  \end{enumerate}
\end{itemize}
\end{frame}

\begin{frame}{各章の概要}
\protect\hypertarget{ux5404ux7ae0ux306eux6982ux8981}{}
\begin{itemize}
\tightlist
\item
  第1章：イントロダクション（Rの基本操作）
\item
  第2章：無作為化実験による因果効果の推定（説明）
\item
  第3章：社会調査研究による母集団の特徴の推論（測定）
\item
  第4章：線形回帰を用いた結果の予測（予測）
\item
  第5章：観察データによる因果効果の推定（説明）
\item
  第6章：確率
\item
  第7章：不確実性の数値化
\end{itemize}
\end{frame}

\hypertarget{ux306aux305cux30c7ux30fcux30bfux5206ux6790ux3092ux5b66ux3076ux306eux304b}{%
\section{1.4
なぜデータ分析を学ぶのか？}\label{ux306aux305cux30c7ux30fcux30bfux5206ux6790ux3092ux5b66ux3076ux306eux304b}}

\begin{frame}{データ分析スキルの重要性}
\protect\hypertarget{ux30c7ux30fcux30bfux5206ux6790ux30b9ux30adux30ebux306eux91cdux8981ux6027}{}
\begin{itemize}
\tightlist
\item
  社会科学者として、母集団の特徴を測り、予測を行い、因果関係を解明するためにデータが必要。
\item
  意思決定者に情報を提供するデータアナリストの需要が高まっている。
\item
  ビッグデータ時代において、計量研究に精通することは非常に価値のあるスキル。
\end{itemize}
\end{frame}

\begin{frame}{コードの学習}
\protect\hypertarget{ux30b3ux30fcux30c9ux306eux5b66ux7fd2}{}
\begin{itemize}
\tightlist
\item
  コンピュータに特定の命令（プログラミング言語）を与えることで分析を実行する。
\item
  本書では、世界中で広く利用されている統計プログラミング言語 \textbf{R}
  を使用する。
\end{itemize}
\end{frame}

\hypertarget{ux6e96ux5099-rux3068rstudio}{%
\section{1.5 準備 (RとRStudio)}\label{ux6e96ux5099-rux3068rstudio}}

\begin{frame}{RとRStudio}
\protect\hypertarget{rux3068rstudio}{}
\begin{itemize}
\tightlist
\item
  \textbf{R}: 統計計算とグラフィックスのためのソフトウェア（エンジン）。
\item
  \textbf{RStudio}: Rをより使いやすくするためのインターフェース（IDE）。
\end{itemize}
\end{frame}

\begin{frame}[fragile]{作業ディレクトリの設定}
\protect\hypertarget{ux4f5cux696dux30c7ux30a3ux30ecux30afux30c8ux30eaux306eux8a2dux5b9a}{}
\begin{itemize}
\tightlist
\item
  データセットを読み込む前に、Rにファイルの場所を教える必要がある。
\end{itemize}

\begin{Shaded}
\begin{Highlighting}[]
\CommentTok{\# 作業ディレクトリの設定例}
\CommentTok{\# setwd("\textasciitilde{}/Desktop/DSS") }
\end{Highlighting}
\end{Shaded}
\end{frame}

\hypertarget{rux306eux7d39ux4ecb}{%
\section{1.6 Rの紹介}\label{rux306eux7d39ux4ecb}}

\begin{frame}[fragile]{Rで計算をする}
\protect\hypertarget{rux3067ux8a08ux7b97ux3092ux3059ux308b}{}
\begin{itemize}
\tightlist
\item
  Rを計算機として使用できる。
\end{itemize}

\begin{Shaded}
\begin{Highlighting}[]
\DecValTok{1} \SpecialCharTok{+} \DecValTok{3}
\end{Highlighting}
\end{Shaded}

\begin{verbatim}
## [1] 4
\end{verbatim}
\end{frame}

\begin{frame}[fragile]{オブジェクトの作成}
\protect\hypertarget{ux30aaux30d6ux30b8ux30a7ux30afux30c8ux306eux4f5cux6210}{}
\begin{itemize}
\tightlist
\item
  情報を「オブジェクト」として保存する（\texttt{\textless{}-}
  割り当て演算子）。
\end{itemize}

\begin{Shaded}
\begin{Highlighting}[]
\NormalTok{four }\OtherTok{\textless{}{-}} \DecValTok{1} \SpecialCharTok{+} \DecValTok{3}
\NormalTok{four}
\end{Highlighting}
\end{Shaded}

\begin{verbatim}
## [1] 4
\end{verbatim}

\begin{Shaded}
\begin{Highlighting}[]
\NormalTok{hello }\OtherTok{\textless{}{-}} \StringTok{"hi, nice to meet you"}
\NormalTok{hello}
\end{Highlighting}
\end{Shaded}

\begin{verbatim}
## [1] "hi, nice to meet you"
\end{verbatim}
\end{frame}

\begin{frame}[fragile]{関数の使用}
\protect\hypertarget{ux95a2ux6570ux306eux4f7fux7528}{}
\begin{itemize}
\tightlist
\item
  関数は特定の操作を実行する（\texttt{関数名(引数)}）。
\end{itemize}

\begin{Shaded}
\begin{Highlighting}[]
\FunctionTok{sqrt}\NormalTok{(}\DecValTok{16}\NormalTok{)}
\end{Highlighting}
\end{Shaded}

\begin{verbatim}
## [1] 4
\end{verbatim}

\begin{Shaded}
\begin{Highlighting}[]
\FunctionTok{sqrt}\NormalTok{(four)}
\end{Highlighting}
\end{Shaded}

\begin{verbatim}
## [1] 2
\end{verbatim}
\end{frame}

\hypertarget{ux30c7ux30fcux30bfux3092ux8aadux307fux8fbcux3093ux3067ux7406ux89e3ux3059ux308b}{%
\section{1.7
データを読み込んで理解する}\label{ux30c7ux30fcux30bfux3092ux8aadux307fux8fbcux3093ux3067ux7406ux89e3ux3059ux308b}}

\begin{frame}[fragile]{データの読み込み}
\protect\hypertarget{ux30c7ux30fcux30bfux306eux8aadux307fux8fbcux307f}{}
\begin{itemize}
\tightlist
\item
  \texttt{read.csv()}
  関数を用いてCSVファイルを読み込み、オブジェクトに格納する。
\end{itemize}

\begin{Shaded}
\begin{Highlighting}[]
\CommentTok{\# データセットの読み込み}
\NormalTok{star }\OtherTok{\textless{}{-}} \FunctionTok{read.csv}\NormalTok{(}\StringTok{"STAR.csv"}\NormalTok{)}
\end{Highlighting}
\end{Shaded}
\end{frame}

\begin{frame}[fragile]{データの理解}
\protect\hypertarget{ux30c7ux30fcux30bfux306eux7406ux89e3}{}
\begin{itemize}
\tightlist
\item
  読み込んだデータの中身や構造を確認する。
\end{itemize}

\begin{Shaded}
\begin{Highlighting}[]
\CommentTok{\# 最初の6行を表示}
\FunctionTok{head}\NormalTok{(star)}
\end{Highlighting}
\end{Shaded}

\begin{verbatim}
##   classtype reading math graduated
## 1     small     578  610         1
## 2   regular     612  612         1
## 3   regular     583  606         1
## 4     small     661  648         1
## 5     small     614  636         1
## 6   regular     610  603         0
\end{verbatim}

\begin{Shaded}
\begin{Highlighting}[]
\CommentTok{\# 次元の確認（行数, 列数）}
\FunctionTok{dim}\NormalTok{(star)}
\end{Highlighting}
\end{Shaded}

\begin{verbatim}
## [1] 1274    4
\end{verbatim}
\end{frame}

\begin{frame}{変数の種類}
\protect\hypertarget{ux5909ux6570ux306eux7a2eux985e}{}
\begin{itemize}
\tightlist
\item
  \textbf{文字変数} (Character variables): テキストデータ。
\item
  \textbf{数値変数} (Numeric variables): 数字。

  \begin{itemize}
  \tightlist
  \item
    \textbf{二値変数} (Binary variables):
    0か1の2つの値のみ（ダミー変数）。
  \item
    \textbf{非二値変数} (Non-binary variables):
    それ以外の数値（例：テストの点数）。
  \end{itemize}
\end{itemize}
\end{frame}

\hypertarget{ux5e73ux5747ux5024ux3092ux8a08ux7b97ux3057ux89e3ux91c8ux3059ux308b}{%
\section{1.8
平均値を計算し、解釈する}\label{ux5e73ux5747ux5024ux3092ux8a08ux7b97ux3057ux89e3ux91c8ux3059ux308b}}

\begin{frame}[fragile]{変数へのアクセス}
\protect\hypertarget{ux5909ux6570ux3078ux306eux30a2ux30afux30bbux30b9}{}
\begin{itemize}
\tightlist
\item
  \texttt{\$} 記号を用いて、データフレーム内の特定の変数にアクセスする。
\end{itemize}

\begin{Shaded}
\begin{Highlighting}[]
\CommentTok{\# reading変数の最初の数値を表示}
\FunctionTok{head}\NormalTok{(star}\SpecialCharTok{$}\NormalTok{reading)}
\end{Highlighting}
\end{Shaded}

\begin{verbatim}
## [1] 578 612 583 661 614 610
\end{verbatim}
\end{frame}

\begin{frame}[fragile]{平均値の計算}
\protect\hypertarget{ux5e73ux5747ux5024ux306eux8a08ux7b97}{}
\begin{itemize}
\tightlist
\item
  \texttt{mean()} 関数を用いて平均値を計算する。
\end{itemize}

\begin{Shaded}
\begin{Highlighting}[]
\CommentTok{\# 読解力テストの平均点}
\FunctionTok{mean}\NormalTok{(star}\SpecialCharTok{$}\NormalTok{reading)}
\end{Highlighting}
\end{Shaded}

\begin{verbatim}
## [1] 628.803
\end{verbatim}

\begin{Shaded}
\begin{Highlighting}[]
\CommentTok{\# 卒業率（二値変数）の平均}
\FunctionTok{mean}\NormalTok{(star}\SpecialCharTok{$}\NormalTok{graduated)}
\end{Highlighting}
\end{Shaded}

\begin{verbatim}
## [1] 0.8697017
\end{verbatim}
\end{frame}

\begin{frame}[fragile]{解釈}
\protect\hypertarget{ux89e3ux91c8}{}
\begin{itemize}
\tightlist
\item
  二値変数の平均値に100を掛けると、その特徴を持つ割合（\%）になる。
\item
  例：\texttt{graduated}
  の平均が0.87であれば、約87\%の生徒が卒業したことを意味する。
\end{itemize}
\end{frame}

\hypertarget{ux307eux3068ux3081}{%
\section{1.9 まとめ}\label{ux307eux3068ux3081}}

\begin{frame}[fragile]{第1章のまとめ}
\protect\hypertarget{ux7b2c1ux7ae0ux306eux307eux3068ux3081}{}
\begin{itemize}
\tightlist
\item
  RとRStudioの基本的な操作を学んだ。
\item
  データの読み込み方法、データフレームの構造、変数の種類を理解した。
\item
  \texttt{mean()}
  関数を用いて記述統計量（平均値）を算出し、その解釈方法を学んだ。
\item
  次章からは、いよいよ実際の研究例を用いて因果関係の推定（説明）に取り組む。
\end{itemize}
\end{frame}

\end{document}
